\documentclass[10pt,a4paper]{amsart}
\usepackage[margin=19mm]{geometry}
\usepackage{amsmath,amssymb,mathtools}
\usepackage[scr=rsfs,cal=euler]{mathalfa}
\usepackage{xfrac}
\usepackage[unicode,hidelinks,bookmarks=false]{hyperref}
\newcommand{\ie}{i.e.\ }
\newcommand{\cf}{cf.\ }
\newcommand{\resp}{respectively\ }
\newcommand{\Subsection}{\subsection}
\providecommand{\nobreakdash}{\nobreak\hbox{-}}
\setlength{\emergencystretch}{2em}
\multlinegap0pt
\usepackage{amssymb}
\usepackage{mathtools}
\usepackage{xcolor}
\usepackage{tikz}
\newtheorem{proposition}{Proposition}
\newcommand{\R}{\ensuremath{\mathbb{R}}}
\newcommand\steq[1]{\stackrel{\text{\rm #1.}}{=}}
\title{Thermodynamic Limits of Stochastic Chemical Reaction Networks with Phosphorylation}
\author{Lucie Laurence, Philippe Robert}
\date{Source: arXiv:2606.30325v1 (CC BY 4.0)}
\begin{document}
\maketitle

\noindent\emph{Source and licence.} Thermodynamic Limits of Stochastic Chemical Reaction Networks with Phosphorylation. Version arXiv:2606.30325v1 (CC BY 4.0), distributed by arXiv under the Creative Commons Attribution 4.0 International licence, http://creativecommons.org/licenses/by/4.0/, as recorded in the arXiv licence metadata for this exact version and shown on the abstract page https://arxiv.org/abs/2606.30325v1. Full licence notice: \texttt{licenses/arxiv-2606.30325-CC-BY-4.0.txt}.\par\smallskip

\noindent\textit{Editorial scope.} Counted unit: the article's proposition M2ExFP (source0.tex lines 1191--1222). The article's complete original proof of it is delivered with it (source0.tex lines 1223--1352, one \texttt{proof} environment of 130 lines). No mathematics is written, repaired or re-derived: the statement and the proof are the article's own lines, in the article's own notation, and no retained line is substituted or re-flowed.  Retained uncounted as the source-local context the proof needs: lines 1121--1124; lines 1171--1175; lines 1177--1181; lines 1183--1186.  Nothing was omitted from any retained span, so the package carries no omission record.  No retained line contains a citation, so the wrapper carries no bibliography and no bibliography entry.  No retained line contains a figure environment, an image-inclusion command or drawing code, so no image is delivered and the package declares no asset dependency.  Editorial formatting only, in addition: an AMS \texttt{amsart} preamble that restates the article's own declarations and macros used by the retained lines, namely the environments \texttt{proposition} on separate counters, together with the standard packages those lines need.  The retained environments print in this document as Proposition 1 in order of appearance, whereas the article prints the same items with the numbering of its own full document.  The complete native source of the article as distributed in the arXiv e-print for this exact version is bundled unchanged as \texttt{sources/204011/source0.tex}.
\begin{multline}\label{VM2}
\lim_{N\to+\infty}\left(\left(\frac{V_N^{H^c}(t)}{N},t{<}t_0\right),\Lambda_N^{H,t_0}\right)\\=\left(\left({u}_A(t),{a}_1(t),{a}_2(t),{u}_B(t),{b}_1(t),{b}_2(t), t{<}t_0\right),\Lambda_\infty^{t_0}\right)\\
{\steq{def}\left(\left(e{-}{a}_1(t){-}{a}_2(t),{a}_1(t),{a}_2(t),f{+}{a}_1(t){+}{a}_2(t){-}1,{b}_1(t),1{-}{a}_1(t){-}{a}_2(t){-}{b}_1(t), t{<}t_0\right),\Lambda_\infty^{t_0}\right)}
\end{multline}

\begin{equation}\label{abb}
a_2^*=\frac{\mu_1  \left(\mu_2 -(\lambda_1  + \mu_2)a_1   \right)}{\mu_2  \left(\lambda_2  +\mu_1  \right)},
\quad b^*_1=\frac{\lambda_2  \left(\mu_2-(\lambda_1 +\mu_2)a_1 \right)}{\mu_2  \left(\lambda_2  +\mu_1  \right)},\quad
b_2^*=\frac{\lambda_1}{\mu_2} a_1^*,
\end{equation}

\begin{multline}\label{PolyP}
P(x)\steq{def}  (\lambda_1\mu_1{-}\lambda_2\mu_2)(\alpha_2\lambda_1\mu_2 (\lambda_2{+}\mu_1){-}\beta_2\lambda_2\mu_1(\lambda_1{+}\mu_2))x^2\\
{+}\mu_2\left[\rule{0mm}{5mm}\alpha_2\lambda_1\mu_2(\lambda_2{+}\mu_1)((\lambda_2{+}\mu_1)e{-}\mu_1){+}\beta_2\lambda_2\mu_1\left((\lambda_1{+}\mu_2)((\lambda_2{+}\mu_1)f{-}\lambda_2){+}\lambda_1\mu_1{-}\lambda_2\mu_2\right)\rule{0mm}{5mm}\right]x\\
  +\beta_2  \lambda_2  \mu_1 \mu_2^{2} (\lambda_2{-}(\lambda_2{+}\mu_1)f),
\end{multline}

\begin{equation}\label{sumC}
a_1^*+a_2^*=\frac{\mu_1\mu_2+a_1^*(\lambda_2\mu_2-\lambda_1\mu_1)}{\mu_2 (\lambda_2  +\mu_1)},\quad
b_1^*+b_2^*=\frac{\lambda_2  \mu_2+a_1^*(\lambda_1\mu_1-\lambda_2\mu_2)}{\mu_2  \left(\lambda_2  +\mu_1  \right)}.
\end{equation}

\begin{proposition}[Equilibrium of the dynamical system of Relation~\eqref{VM2}]\label{M2ExFP}\ \\
Under the condition $e{+}f>1$,
\begin{enumerate}
\item if $\lambda_1\mu_1{-}\lambda_2\mu_2>0$.

\medskip
\noindent  If $e{\not=}\mu_2/(\lambda_1{+}\mu_2)$ and $f{\not=}\lambda_2/(\lambda_2{+}\mu_1)$,
there exists an equilibrium point if and only if the conditions
\begin{equation}\label{Cond21}
e > \frac{\mu_2}{\lambda_1{+}\mu_2} \text{ and }   f >\frac{\lambda_2}{\lambda_2{+}\mu_1},
\end{equation}
hold and in this case, it is unique. 
\item If $\lambda_1\mu_1{-}\lambda_2\mu_2<0$.\medskip
  \begin{enumerate}
  \item If the conditions
\begin{equation}\label{Cond2}
e > \frac{\mu_1}{\lambda_2{+}\mu_1} \text{ and }   f >\frac{\lambda_1}{\lambda_1{+}\mu_2},
\end{equation}
and 
\begin{equation}\label{Cond3}
\left(e{-}\frac{\mu_2}{\lambda_1{+}\mu_2}\right)\left(f {-}\frac{\lambda_2}{\lambda_2{+}\mu_1}\right) > 0,
\end{equation}
hold, then there exists a unique equilibrium point.
\item If Condition~\eqref{Cond2}
and relation
\begin{equation}\label{Cond32}
\left(e{-}\frac{\mu_2}{\lambda_1{+}\mu_2}\right)\left(f {-}\frac{\lambda_2}{\lambda_2{+}\mu_1}\right) < 0,
\end{equation}
hold, then there are  either $0$ or $2$ equilibrium points: There exists a subset $S$, resp. $S_0$, of $\R_+^2$ with non-empty interior such that if $(\alpha_2,\beta_2){\in}S$, resp. $(\alpha_2,\beta_2){\in}S_0$, then there are two equilibrium points, resp. there does not exist an equilibrium point. 
\end{enumerate}
\end{enumerate}
\end{proposition}

\begin{proof}
With Relation~\eqref{abb}, a possible equilibrium point $(a_1^*,a_2^*,b_1^*,b_2^*)$ should satisfy the relations
\begin{equation}\label{feq1}
0{<} a_1^*{<} C_0\steq{def}\frac{\mu_2}{\lambda_1  + \mu_2}, \quad a_1^*{+}a_2^*{<}e \text{ and } b_1^*{+}b_2^*{<}f,
\end{equation}
and therefore, with Relations~\eqref{sumC}, the inequalities
\begin{equation}\label{feq12}
\frac{\mu_1}{\lambda_2  + \mu_1}{-}e < a_1^*\frac{\lambda_1\mu_1{-}\lambda_2\mu_2}{\mu_2(\lambda_2  + \mu_1)}
< f{-} \frac{\lambda_2}{\lambda_2  +\mu_1}.
\end{equation}
As it can be seen this implies that the relation $e{+}f{>}1$ holds. 
%\[\mu_2(\mu_1-(\lambda_2  + \mu_1)e)\le a_1^*(\lambda_1\mu_1-\lambda_2\mu_2)\le \mu_2((\lambda_2  + \mu_1)f-\lambda_2).\]
%This  implies in particular that $e{+}f{\ge}1$.

We define
\[
C_1\steq{def}\mu_2\frac{\mu_1{-}(\lambda_2{+}\mu_1)e}{\lambda_1\mu_1{-}\lambda_2\mu_2} \text{ and } C_2\steq{def}\mu_2\frac{(\lambda_2{+}\mu_1)f{-} \lambda_2}{\lambda_1\mu_1{-}\lambda_2\mu_2},
\]
note that
\begin{equation}\label{Ineq}
\begin{cases}
&C_2{-}C_1{=}\displaystyle\frac{\mu_2(\lambda_2{+}\mu_1)(e{+}f{-}1)}{\lambda_1\mu_1{-}\lambda_2\mu_2}, \quad C_2{-}C_0{=}\frac{\mu_2  \left(\lambda_2  +\mu_1  \right) \left((\lambda_1{+}\mu_2)f{-} \lambda_1 \right)}{\left(\mu_2  +\lambda_1  \right) \left(\lambda_1  \mu_1{-}\lambda_2  \mu_2  \right)},\\
\ \\
 &C_0{-}C_1{=}\displaystyle\frac{\mu_2(\lambda_2{+}\mu_1)((\lambda_1{+}\mu_2)e{-}\mu_2)}{(\lambda_1\mu_1{-}\lambda_2\mu_2)(\mu_2{+}\lambda_1)}.
\end{cases}
\end{equation}
With Relation~\eqref{PolyP} (and trite calculations), we obtain
\begin{equation}\label{RelP}
\begin{cases}
&P(0)\displaystyle=\beta_2  \lambda_2  \mu_1 \mu_2^{2} (\lambda_2{-}(\lambda_2{+}\mu_1)f )\smallskip\\
&P\left(C_0\right)\displaystyle=\mu_2^{3} \alpha_2  \lambda_1  \left(\lambda_2 {+}\mu_1  \right)^{2} \left((\lambda_1{+}\mu_2)e  {-}\mu_2  \right)\smallskip\\
&P(C_1)\displaystyle=\mu_1 \mu_2^{2} \beta_2  \lambda_2  \left(\lambda_2{+}\mu_1  \right)^{2} \left(\lambda_1  \mu_1  {-}\lambda_2  \mu_2  \right) \left(\mu_2{-}(\lambda_1{+} \mu_2)e\right) \left(e{+}f{-}1\right)\smallskip\\
&P(C_2)\displaystyle=\mu_2^{3} \alpha_2  \lambda_1\left(\lambda_1  \mu_1  {-}\lambda_2  \mu_2\right)  \left(\lambda_2 {+}\mu_1  \right)^{2} \left((\lambda_2{+}\mu_1)f  {-}\lambda_2  \right) \left(e{+}f{-}1\right).
\end{cases}
\end{equation}

If $\lambda_1\mu_1-\lambda_2\mu_2>0$ holds, to have an equilibrium $(a_1^*,a_2^*,b_1^*,b_2^*)$ for the dynamical system, the root $a_1^*$ of the polynomial $P$ must be  such that  $a_1^*{\in}(C_1{\vee}0,C_2{\wedge}C_0)$ because of Relation~\eqref{feq12}.  In particular, we must have the relations $C_2{>}0$ and $C_1{<}C_0$, which imply, by using Inequalities~\eqref{Ineq}, that the relations
\[
f> \frac{\lambda_2}{\lambda_2{+}\mu_1} \text{ and } e >\frac{\mu_2}{\lambda_1{+}\mu_2}
\]
hold. With Relations~\eqref{RelP}, we obtain the inequalities  $P(C_1){<}0$  and $P(C_0){>}0$ therefore there is a unique equilibrium point $(a_1^*,a_2^*,b_1^*,b_2^*)$ in $\R_+^4$ with $a_1^*{\in}(C_1^+,C_0{\wedge}C_2)$. The first case is proved. 


For the second part of the proposition, when $\lambda_1\mu_1{-}\lambda_2\mu_2{<}0$, we must have $a_1^*{\in}(C_2^+,C_1{\wedge}C_0)$ and therefore $C_1{>}0$ and $C_2{<}C_0$, and hence the conditions
\[
e>\frac{\mu_1}{\lambda_2{+}\mu_1} \text{ and }f>\frac{\lambda_1}{\lambda_1{+}\mu_2}.
\]

From now on, assume that Condition~\eqref{Cond2} holds.

If
\[
e> \frac{\mu_2}{\lambda_1{+}\mu_2} \text{ and } f> \frac{\lambda_2}{\lambda_2{+}\mu_1},
\]
then  $C_2{<}0{<}C_0{<}C_1$ and by using again Relations~\eqref{RelP}, we have $P(0){<}0$ and $P(C_0){>}0$. Hence, there is a unique equilibrium point $(a_1^*,a_2^*,b_1^*,b_2^*)$ in $\R_+^4$ with $a_1^*{\in}(0,C_1)$. 

Otherwise if
\[
e< \frac{\mu_2}{\lambda_1{+}\mu_2} \text{ and } f< \frac{\lambda_2}{\lambda_2{+}\mu_1},
\]
then,  $0{<}C_2{<}C_1{<}C_0$ and $P(C_2){>}0$ and $P(C_1){<}0$. Hence, there is a unique equilibrium point $(a_1^*,a_2^*,b_1^*,b_2^*)$ in $\R_+^4$ with $a_1^*{\in}(C_2,C_1)$. 

For the proof of the last part of the proposition, we introduce $W_0$ the location of the extremum of $P$,
 \[
 W_0=-\frac{\mu_2\left[\rule{0mm}{5mm}\alpha_2\lambda_1\mu_2(\lambda_2{+}\mu_1)((\lambda_2{+}\mu_1)e{-}\mu_1){+}\beta_2\lambda_2\mu_1\left(\rule{0mm}{4mm}(\lambda_1{+}\mu_2)((\lambda_2{+}\mu_1)f{-}\lambda_2){+}\lambda_1\mu_1{-}\lambda_2\mu_2\right)\rule{0mm}{5mm}\right]}{2(\lambda_1\mu_1{-}\lambda_2\mu_2)(\alpha_2\lambda_1\mu_2 (\lambda_2{+}\mu_1){-}\beta_2\lambda_2\mu_1(\lambda_1{+}\mu_2))}.
 \]
With some calculations, we have
 \[
 C_0{-}W_0=\frac{\alpha_2\lambda_1\mu_2((\lambda_2{+}\mu_1)((\lambda_1{+}\mu_2)e{-}\mu_2)+\lambda_1\mu_1{-}\lambda_2\mu_2)+ \beta_2\lambda_2\mu_1(\mu_2{+}\lambda_1)((\lambda_1{+}\mu_2)f{-}\lambda_1)}{2(\lambda_1\mu_1{-}\lambda_2\mu_2)(\alpha_2\lambda_1\mu_2 (\lambda_2{+}\mu_1){-}\beta_2\lambda_2\mu_1(\lambda_1{+}\mu_2))}
 \]
 and
  \[
  W_0{-}C_2= -(\lambda_2{+}\mu_1)\mu_2
  \frac{\alpha_2((\lambda_2{+}\mu_1)(e{+}f{-}1){+}(\lambda_2{+}\mu_1)f{-}\lambda_2){+}\beta_2\lambda_2\mu_1(\lambda_1{-}(\lambda_1{+}\mu_2)f)}{2(\lambda_1\mu_1{-}\lambda_2\mu_2)(\alpha_2\lambda_1\mu_2 (\lambda_2{+}\mu_1){-}\beta_2\lambda_2\mu_1(\lambda_1{+}\mu_2))}
  \]

If Condition~\eqref{Cond2} and
\[
e>\frac{\mu_2}{\lambda_1{+}\mu_2} \text{ and } f< \frac{\lambda_2}{\lambda_2{+}\mu_1}
\]
hold, we can write
\[
f=w\frac{\lambda_1}{\lambda_1+\mu_2} +(1-w)\frac{\lambda_2}{\lambda_2{+}\mu_1},
\]
with $w{\in}(0,1)$. We have $0{<}C_2{<}C_0{<}C_1$ and $P(C_2){>}0$ and $P(C_0){>}0$. Hence there are either $0$ or $2$ roots (counting the multiplicity if $P(W_0){=}0$), of $P$ in $(C_2,C_0)$. 

It is easily checked that 
\[
\lim_{(\alpha_2,\beta_2)\to(0,1)} (W_0{-}C_2,C_0{-}W_0)=(1{-}w)\frac{\mu_2}{2(\mu_2{+}\lambda_1)}(1,1)
\]
and, with some calculations, 
\[
\lim_{(\alpha_2,\beta_2)\to(0,1)} P(W_0)=\mu_2^2\mu_1^3\lambda_2^3(1{-}w)^2(\mu_2{+}\lambda_1)^2(\lambda_1\mu_1{-}\lambda_2\mu_2)^3 <0. 
\]

Hence there exists a neighborhood $V$ of $(1,0)$ such that if $(\alpha_2,\beta_2){\in}V$, then $W_0{\in}(C_2,C_0)$ and $P(W_0)<0$. In this case there are two roots of $P$ in $(C_2,C_0)$.

The result
\[
\lim_{(\alpha_2,\beta_2)\to(1,0)} P(W_0)=\mu_2^5\lambda_1^3(\lambda_2{+}\mu_1)^3(\lambda_1{+}\mu_2)(\lambda_2\mu_2{-}\lambda_1\mu1)((\lambda_2{+}\mu_1)e{-}\mu_1)^2>0,
\]
shows that, for $(\alpha_2,\beta_2)$ in a small neighborhood of $(1,0)$: if $W_0{\in}(C_2,C_1)$ then there does not exist a root of $P$ in $(C_2,C_1)$, otherwise, if $W_0{\not\in}(C_2,C_1)$ the same conclusion holds obviously.

If Condition~\eqref{Cond2} and
\[
e<\frac{\mu_2}{\lambda_1{+}\mu_2} \text{ and } f>\frac{\lambda_2}{\lambda_2{+}\mu_1}
\]
we can write
\[
e=w\frac{\mu_2}{\lambda_1+\mu_2} +(1-w)\frac{\mu_1}{\lambda_2{+}\mu_1},
\]
with $w{\in}(0,1)$. We have $C_2{<}0{<}C_1{<}C_0$ and $P(0){<}0$ and $P(C_1){<}0$. There is either $0$ or $2$ roots of $P$ in $(0,C_1)$. 

\[
\lim_{(\alpha_2,\beta_2)\to(1,0)} (W_0,C_1{-}W_0)=w\frac{\mu_2}{2(\mu_2{+}\lambda_1)}(1,1)
\]
and
\[
\lim_{(\alpha_2,\beta_2)\to(1,0)} P(W_0)=(\lambda_2{+}\mu_1)^3(\lambda_2\mu_2{-}\lambda_1\mu_1)^3\mu_2^5\lambda_1^3w^2>0,
\]
there are two roots of $P$ in $(0,C_1)$ for $(\alpha_2,\beta_2)$ in a small neighborhood of $(1,0)$.

Finally, the relation
\[
\lim_{(\alpha_2,\beta_2)\to(0,1)} P(W_0)=\lambda_2^3\mu_1^3\mu_2^2(\lambda_2{+}\mu_1)^2(\lambda_1{+}\mu_2)^3(\lambda_1\mu_1{-}\lambda_2\mu_2)((\lambda_1{+}\mu_2)f{-}\lambda_1)^2<0,
\]
shows that there does not exist a root of $P$ in $(0,C_1)$ for $(\alpha_2,\beta_2)$ in a small neighborhood of $(0,1)$.

The proposition is proved.
\end{proof}

\end{document}
